\documentclass[aspectratio=169]{beamer}
\usetheme{Madrid}
\usecolortheme{default}
\usepackage[utf8]{inputenc}
\usepackage[spanish,es-nodecimaldot]{babel}
\usepackage{amsmath}
\usepackage{amssymb}
\usepackage{booktabs}
\usepackage{multirow}
\usepackage{graphicx}

\definecolor{unsblue}{RGB}{30,60,120}
\setbeamercolor{structure}{fg=unsblue}
\setbeamercolor{palette primary}{bg=unsblue,fg=white}
\setbeamercolor{palette secondary}{bg=unsblue!80,fg=white}
\setbeamercolor{palette tertiary}{bg=unsblue!60,fg=white}

\setbeamertemplate{navigation symbols}{}
\setbeamertemplate{footline}[frame number]

\title[MAD1 - Clase 27]{Métodos de Análisis de Datos 1}
\subtitle{Clase 27: Modelos en Producción e Interpretabilidad}
\author{Prof. José}
\institute{Departamento de Matemática \\ Universidad Nacional del Sur}
\date{}

\begin{document}

\begin{frame}
\titlepage
\end{frame}

\begin{frame}{Contenidos de la Clase}
\tableofcontents
\end{frame}

\section{Introducción}

\begin{frame}{¿Por qué esta unidad?}
Hasta ahora aprendimos a:
\begin{itemize}
\item Obtener y limpiar datos
\item Explorar y visualizar
\item Construir y evaluar modelos
\end{itemize}

\vspace{0.3cm}

\textbf{Pero falta:}
\begin{itemize}
\item ¿Cómo llevamos modelos a producción?
\item ¿Cómo explicamos sus decisiones?
\item ¿Cómo comunicamos resultados?
\item ¿Qué responsabilidades éticas tenemos?
\end{itemize}

\vspace{0.3cm}

Estos aspectos son \textbf{fundamentales} para el éxito profesional en ciencia de datos.
\end{frame}

\section{Ciclo de Vida de Modelos}

\begin{frame}{Fases del Ciclo de Vida}
\begin{enumerate}
\item \textbf{Desarrollo:} exploración, experimentación, validación
\item \textbf{Deployment:} poner el modelo en producción
\item \textbf{Monitoreo:} supervisar rendimiento y detectar problemas
\item \textbf{Reentrenamiento:} actualizar con nuevos datos
\end{enumerate}

\vspace{0.3cm}

\begin{block}{Realidad}
El desarrollo del modelo es solo el 20\% del trabajo. El 80\% restante es deployment, mantenimiento y monitoreo.
\end{block}
\end{frame}

\begin{frame}{1. Desarrollo e Investigación}
\textbf{Actividades:}
\begin{itemize}
\item Análisis exploratorio de datos (EDA)
\item Ingeniería de features
\item Experimentación con algoritmos
\item Validación rigurosa
\item Análisis de errores
\end{itemize}

\vspace{0.3cm}

\textbf{Herramientas:}
\begin{itemize}
\item Jupyter notebooks, VS Code
\item scikit-learn, statsmodels
\item pandas, matplotlib, seaborn
\end{itemize}

\vspace{0.3cm}

\textbf{Output:} Modelo entrenado que cumple requisitos de performance y negocio
\end{frame}

\begin{frame}{2. Deployment (Puesta en Producción)}
\textbf{Objetivo:} Hacer el modelo accesible para usuarios o sistemas

\vspace{0.3cm}

\textbf{Pasos típicos:}
\begin{enumerate}
\item \textbf{Serialización:} guardar modelo entrenado
\item \textbf{API REST:} exponer modelo vía endpoints HTTP
\item \textbf{Containerización:} empaquetar con Docker
\item \textbf{Integración:} conectar con bases de datos, pipelines
\item \textbf{Deploy:} subir a servidor/cloud (AWS, Azure, GCP)
\end{enumerate}

\vspace{0.3cm}

\begin{exampleblock}{Ejemplo de flujo}
Notebook → Serializar → Flask/FastAPI → Docker → AWS/Azure
\end{exampleblock}
\end{frame}

\begin{frame}{3. Monitoreo}
\textbf{¿Qué monitorear?}

\vspace{0.2cm}

\textbf{Performance del modelo:}
\begin{itemize}
\item Accuracy, precision, recall, AUC
\item Distribución de predicciones
\item Comparar predicciones vs valores reales
\end{itemize}

\vspace{0.2cm}

\textbf{Calidad de datos:}
\begin{itemize}
\item Distribución de features de entrada
\item Valores faltantes, outliers
\item Data drift (cambios en $P(X)$)
\end{itemize}

\vspace{0.2cm}

\textbf{Aspectos técnicos:}
\begin{itemize}
\item Latencia (tiempo de respuesta)
\item Tasa de errores
\item Disponibilidad del servicio
\end{itemize}
\end{frame}

\begin{frame}{Data Drift vs Concept Drift}
\begin{columns}
\begin{column}{0.5\textwidth}
\textbf{Data Drift}

Cambio en la distribución de las entradas: $P(X)$

\vspace{0.2cm}

\textbf{Ejemplos:}
\begin{itemize}
\item Perfiles de clientes cambian post-pandemia
\item Variables demográficas evolucionan
\item Estacionalidad en ventas
\end{itemize}

\vspace{0.2cm}

\textbf{Detección:}
\begin{itemize}
\item Test Kolmogorov-Smirnov
\item Comparar estadísticos (media, varianza)
\item Histogramas comparativos
\end{itemize}
\end{column}

\begin{column}{0.5\textwidth}
\textbf{Concept Drift}

Cambio en la relación: $P(Y|X)$

\vspace{0.2cm}

\textbf{Ejemplos:}
\begin{itemize}
\item Regulaciones cambian
\item Estafadores adaptan métodos
\item Preferencias de usuarios evolucionan
\end{itemize}

\vspace{0.2cm}

\textbf{Detección:}
\begin{itemize}
\item Monitorear accuracy en producción
\item Comparar predicciones vs reales
\item Alertas cuando cae performance
\end{itemize}
\end{column}
\end{columns}
\end{frame}

\begin{frame}{4. Reentrenamiento}
\textbf{¿Cuándo reentrenar?}

\begin{itemize}
\item \textbf{Programado:} automático cada cierto tiempo (diario, semanal, mensual)
\item \textbf{Disparado:} cuando métricas caen bajo umbral
\item \textbf{Manual:} cuando hay cambios importantes conocidos
\end{itemize}

\vspace{0.3cm}

\textbf{Buenas prácticas:}
\begin{itemize}
\item \textbf{Versionado:} mantener registro de todas las versiones
\item \textbf{A/B testing:} comparar nuevo vs antiguo antes de reemplazar
\item \textbf{Rollback:} poder volver a versión anterior si falla
\item \textbf{Documentación:} registrar cambios, datos usados, métricas
\end{itemize}

\vspace{0.3cm}

\textbf{Herramientas:} MLflow, DVC, Neptune.ai, Weights \& Biases
\end{frame}

\begin{frame}{Consideraciones de Deployment}
\textbf{Escalabilidad:}
\begin{itemize}
\item Latencia: tiempo de respuesta (usualmente $<$ 100ms)
\item Throughput: predicciones por segundo
\item Batch vs real-time
\end{itemize}

\vspace{0.3cm}

\textbf{Mantenibilidad:}
\begin{itemize}
\item Documentación completa
\item Tests unitarios
\item Logging de predicciones y errores
\item Reproducibilidad (código, datos, modelo)
\end{itemize}

\vspace{0.3cm}

\textbf{Trade-offs:}
\begin{itemize}
\item Modelos complejos: más precisos pero más lentos
\item Simplicidad vs performance
\item Costo computacional vs mejora marginal
\end{itemize}
\end{frame}

\section{Serialización de Modelos}

\begin{frame}{¿Qué es Serialización?}
\begin{block}{Definición}
Convertir un objeto Python (modelo entrenado) en un formato que puede guardarse en disco y recargarse posteriormente.
\end{block}

\vspace{0.3cm}

\textbf{¿Por qué serializar?}
\begin{itemize}
\item No reentrenar cada vez que se usa el modelo
\item Compartir modelos entre equipos
\item Versionado y control de cambios
\item Deployment en producción
\item Backup y recuperación
\end{itemize}
\end{frame}

\begin{frame}[fragile]{pickle: Formato Estándar}
\textbf{Nativo de Python, simple de usar}

\vspace{0.2cm}

\begin{columns}
\begin{column}{0.5\textwidth}
\textbf{Ventajas:}
\begin{itemize}
\item No requiere librerías extra
\item Serializa casi cualquier objeto
\item API simple
\end{itemize}
\end{column}

\begin{column}{0.5\textwidth}
\textbf{Desventajas:}
\begin{itemize}
\item Solo Python
\item Vulnerabilidades de seguridad
\item Incompatibilidades entre versiones
\end{itemize}
\end{column}
\end{columns}

\vspace{0.3cm}

\textbf{Uso básico:}
\begin{verbatim}
import pickle

# Guardar
with open('modelo.pkl', 'wb') as f:
    pickle.dump(modelo, f)

# Cargar
with open('modelo.pkl', 'rb') as f:
    modelo = pickle.load(f)
\end{verbatim}
\end{frame}

\begin{frame}[fragile]{joblib: Optimizado para ML}
\textbf{Mejor para modelos grandes con arrays NumPy}

\vspace{0.2cm}

\textbf{Ventajas:}
\begin{itemize}
\item Más eficiente que pickle
\item Compresión automática
\item API compatible con pickle
\end{itemize}

\vspace{0.3cm}

\textbf{Uso:}
\begin{verbatim}
from joblib import dump, load

# Guardar
dump(modelo, 'modelo.joblib')

# Cargar
modelo = load('modelo.joblib')
\end{verbatim}

\vspace{0.2cm}

\begin{block}{Recomendación}
joblib es la opción preferida para modelos de scikit-learn.
\end{block}
\end{frame}

\begin{frame}[fragile]{ONNX: Formato Interoperable}
\textbf{Open Neural Network Exchange}

\vspace{0.2cm}

\textbf{Ventajas:}
\begin{itemize}
\item Portable entre frameworks (PyTorch, TensorFlow, scikit-learn)
\item Portable entre lenguajes (Python, C++, Java, JavaScript)
\item Optimizado para inferencia rápida
\end{itemize}

\vspace{0.2cm}

\textbf{Uso con scikit-learn:}
\begin{verbatim}
from skl2onnx import convert_sklearn
from skl2onnx.common.data_types import FloatTensorType

initial_type = [('input', FloatTensorType([None, n_features]))]
onx = convert_sklearn(modelo, initial_types=initial_type)

with open("modelo.onnx", "wb") as f:
    f.write(onx.SerializeToString())
\end{verbatim}
\end{frame}

\begin{frame}[fragile]{Buenas Prácticas de Serialización}
\textbf{1. Guardar todo el pipeline:}
\begin{verbatim}
from sklearn.pipeline import Pipeline

pipeline = Pipeline([
    ('scaler', StandardScaler()),
    ('model', RandomForestClassifier())
])
pipeline.fit(X_train, y_train)
dump(pipeline, 'pipeline.joblib')
\end{verbatim}

\vspace{0.3cm}

\textbf{2. Incluir metadata:}
\begin{itemize}
\item Versión del modelo
\item Fecha de entrenamiento
\item Métricas de performance
\item Features utilizadas
\item Versiones de librerías
\end{itemize}
\end{frame}

\begin{frame}[fragile]{Ejemplo Completo de Serialización}
\small
\begin{verbatim}
from sklearn.pipeline import Pipeline
from joblib import dump
import json
from datetime import datetime

# Pipeline completo
pipeline = Pipeline([
    ('scaler', StandardScaler()),
    ('classifier', RandomForestClassifier(n_estimators=100))
])
pipeline.fit(X_train, y_train)

# Guardar pipeline
dump(pipeline, 'pipeline_v1.0.0.joblib')

# Guardar metadata
metadata = {
    'version': '1.0.0',
    'timestamp': datetime.now().isoformat(),
    'accuracy': pipeline.score(X_test, y_test),
    'features': list(X_train.columns)
}
with open('pipeline_v1.0.0_metadata.json', 'w') as f:
    json.dump(metadata, f, indent=2)
\end{verbatim}
\end{frame}

\section{Interpretabilidad de Modelos}

\begin{frame}{¿Por qué Interpretabilidad?}
\textbf{Razones importantes:}

\begin{enumerate}
\item \textbf{Confianza:} usuarios confían más en modelos explicables
\item \textbf{Debugging:} identificar errores y sesgos
\item \textbf{Cumplimiento:} regulaciones (finanzas, salud) requieren explicaciones
\item \textbf{Mejora:} entender qué aprende el modelo ayuda a mejorarlo
\item \textbf{Aceptación:} stakeholders adoptan más fácilmente modelos comprensibles
\end{enumerate}

\vspace{0.3cm}

\begin{alertblock}{Pregunta clave}
¿Puedo explicar a un cliente por qué el modelo tomó esta decisión?
\end{alertblock}
\end{frame}

\begin{frame}{Modelos Interpretables vs Caja Negra}
\begin{columns}
\begin{column}{0.5\textwidth}
\textbf{Inherentemente Interpretables:}

\vspace{0.2cm}

\textbf{Regresión Lineal/Logística:}
\begin{itemize}
\item Coeficientes claros
\item Dirección y magnitud del efecto
\end{itemize}

\vspace{0.2cm}

\textbf{Árboles de Decisión:}
\begin{itemize}
\item Reglas if-then visuales
\item Fácil explicar a no técnicos
\end{itemize}
\end{column}

\begin{column}{0.5\textwidth}
\textbf{Caja Negra:}

\vspace{0.2cm}

\textbf{Random Forest, Boosting:}
\begin{itemize}
\item Múltiples árboles combinados
\item No hay regla simple
\end{itemize}

\vspace{0.2cm}

\textbf{Redes Neuronales:}
\begin{itemize}
\item Millones de parámetros
\item Transformaciones complejas
\end{itemize}

\vspace{0.2cm}

\textbf{SVM con kernels:}
\begin{itemize}
\item Espacios de alta dimensión
\item No directamente interpretables
\end{itemize}
\end{column}
\end{columns}
\end{frame}

\begin{frame}{Trade-off Interpretabilidad vs Performance}
\textbf{Estrategia recomendada:}
\begin{enumerate}
\item Comenzar con modelos simples (baseline)
\item Si performance insuficiente, probar modelos complejos
\item Usar técnicas de interpretabilidad (SHAP, LIME)
\item Documentar trade-offs claramente
\end{enumerate}
\end{frame}

\begin{frame}{SHAP: Valores de Shapley}
\textbf{SHapley Additive exPlanations}

\vspace{0.3cm}

\begin{block}{Idea Principal}
Para cada predicción, SHAP calcula la contribución de cada feature:
$$f(x) = \phi_0 + \sum_{i=1}^{p} \phi_i$$

donde:
\begin{itemize}
\item $f(x)$: predicción para instancia $x$
\item $\phi_0$: valor base (predicción promedio)
\item $\phi_i$: contribución de feature $i$ (valor SHAP)
\end{itemize}
\end{block}

\vspace{0.3cm}

\textbf{Interpretación:}
\begin{itemize}
\item Valor positivo: empuja predicción hacia arriba
\item Valor negativo: empuja predicción hacia abajo
\item Magnitud: importancia de la contribución
\end{itemize}
\end{frame}

\begin{frame}{Propiedades Deseables de SHAP}
\textbf{¿Por qué SHAP es mejor que otros métodos?}

\vspace{0.3cm}

\begin{enumerate}
\item \textbf{Local accuracy:} 
\begin{itemize}
\item La suma de contribuciones reproduce exactamente la predicción
\end{itemize}

\vspace{0.2cm}

\item \textbf{Missingness:}
\begin{itemize}
\item Features ausentes tienen contribución cero
\end{itemize}

\vspace{0.2cm}

\item \textbf{Consistency:}
\begin{itemize}
\item Si un modelo cambia para depender más de una feature, su contribución SHAP aumenta
\end{itemize}
\end{enumerate}

\vspace{0.3cm}

\textbf{Basado en teoría de juegos:} valores de Shapley garantizan reparto justo del pago total entre jugadores.
\end{frame}

\begin{frame}{Gráficos SHAP}
\textbf{1. Force Plot:}
\begin{itemize}
\item Visualiza contribuciones a una predicción individual
\item Muestra cómo cada feature empuja hacia arriba/abajo
\end{itemize}

\vspace{0.2cm}

\textbf{2. Summary Plot:}
\begin{itemize}
\item Distribución de SHAP values para todo el dataset
\item Features ordenadas por importancia
\item Color indica valor alto/bajo de la feature
\end{itemize}

\vspace{0.2cm}

\textbf{3. Dependence Plot:}
\begin{itemize}
\item Relación entre valor de feature y su SHAP value
\item Detecta interacciones
\end{itemize}

\vspace{0.2cm}

\textbf{4. Waterfall Plot:}
\begin{itemize}
\item Descomposición aditiva paso a paso
\end{itemize}
\end{frame}

\begin{frame}[fragile]{SHAP en Python}
\small
\begin{verbatim}
import shap
from sklearn.ensemble import RandomForestClassifier

# Entrenar modelo
modelo = RandomForestClassifier()
modelo.fit(X_train, y_train)

# Crear explainer
explainer = shap.TreeExplainer(modelo)

# Calcular SHAP values
shap_values = explainer.shap_values(X_test)

# Summary plot (importancia global)
shap.summary_plot(shap_values[1], X_test)

# Force plot (predicción individual)
shap.force_plot(explainer.expected_value[1], 
                shap_values[1][0], 
                X_test.iloc[0])
\end{verbatim}
\end{frame}

\begin{frame}{LIME: Explicaciones Locales}
\textbf{Local Interpretable Model-agnostic Explanations}

\vspace{0.3cm}

\textbf{Idea:} Aproximar localmente el modelo complejo con uno simple

\vspace{0.3cm}

\textbf{Procedimiento:}
\begin{enumerate}
\item Generar perturbaciones del punto a explicar
\item Obtener predicciones del modelo original
\item Entrenar modelo simple (regresión lineal) en las perturbaciones
\item El modelo simple aproxima localmente el comportamiento
\end{enumerate}

\vspace{0.3cm}

\textbf{Analogía:} Aproximar función compleja con recta tangente en un punto
\end{frame}

\begin{frame}{LIME: Ventajas y Desventajas}
\begin{columns}
\begin{column}{0.5\textwidth}
\textbf{Ventajas:}
\begin{itemize}
\item Model-agnostic
\item Explicaciones locales
\item Interpretable (usa modelos lineales)
\item Funciona con cualquier tipo de datos
\end{itemize}
\end{column}

\begin{column}{0.5\textwidth}
\textbf{Desventajas:}
\begin{itemize}
\item Inestabilidad
\item Dependencia de hiperparámetros
\item No garantiza consistencia global
\item Puede ser engañoso
\end{itemize}
\end{column}
\end{columns}

\vspace{0.4cm}

\begin{block}{Cuándo usar LIME}
Explicar predicciones individuales específicas (por ejemplo, por qué se rechazó un préstamo a este cliente particular).
\end{block}
\end{frame}

\begin{frame}[fragile]{LIME en Python}
\small
\begin{verbatim}
import lime
import lime.lime_tabular

# Crear explainer
explainer = lime.lime_tabular.LimeTabularExplainer(
    X_train.values,
    feature_names=X_train.columns,
    class_names=['Clase 0', 'Clase 1'],
    mode='classification'
)

# Explicar una instancia
i = 0  # índice de la instancia a explicar
exp = explainer.explain_instance(
    X_test.values[i], 
    modelo.predict_proba,
    num_features=10
)

# Visualizar
exp.show_in_notebook()
exp.as_pyplot_figure()
\end{verbatim}
\end{frame}

\begin{frame}{Partial Dependence Plots (PDP)}
\textbf{Muestran el efecto marginal de una feature sobre la predicción}

\vspace{0.3cm}

\textbf{Cálculo:}
$$PD_{X_j}(x_j) = \frac{1}{n}\sum_{i=1}^{n}f(x_j, x_{-j}^{(i)})$$

Fijamos $X_j$ en un valor, promediamos predicciones sobre todas las demás features.

\vspace{0.3cm}

\textbf{Interpretación:}
\begin{itemize}
\item Eje X: valores de la feature
\item Eje Y: predicción promedio
\item Pendiente positiva: relación positiva
\item Curva plana: no afecta la predicción
\end{itemize}

\vspace{0.3cm}

\textbf{Limitación:} Asume independencia entre features (raramente cierto)
\end{frame}

\begin{frame}[fragile]{PDP en Python}
\begin{verbatim}
from sklearn.inspection import PartialDependenceDisplay

# PDP para una feature
fig, ax = plt.subplots(figsize=(10, 6))
PartialDependenceDisplay.from_estimator(
    modelo, X_train, 
    features=['edad'],
    ax=ax
)

# PDP para múltiples features + interacción
fig, ax = plt.subplots(figsize=(12, 4))
PartialDependenceDisplay.from_estimator(
    modelo, X_train, 
    features=['edad', 'ingreso', ('edad', 'ingreso')],
    ax=ax
)
\end{verbatim}
\end{frame}

\begin{frame}{Feature Importance}
\textbf{Mide importancia relativa de cada feature}

\vspace{0.3cm}

\textbf{Importancia en Árboles (Gini/Entropy):}
\begin{itemize}
\item Incorporada en scikit-learn
\item Rápida de calcular
\item Sesgada hacia features con alta cardinalidad
\end{itemize}

\vspace{0.3cm}

\textbf{Permutation Importance:}
\begin{itemize}
\item Mide aumento en error al permutar una feature
\item Model-agnostic
\item No sesgada por cardinalidad
\item Más lenta de calcular
\item Puede confundir con features correlacionadas
\end{itemize}

\vspace{0.3cm}

\textbf{Recomendación:} Usar permutation importance como verificación de feature importance estándar.
\end{frame}

\begin{frame}[fragile]{Feature Importance en Python}
\small
\begin{verbatim}
# Importancia estándar (árboles)
importances = modelo.feature_importances_
indices = np.argsort(importances)[::-1]

plt.bar(range(len(importances)), importances[indices])
plt.xticks(range(len(importances)), 
           X_train.columns[indices], rotation=90)

# Permutation importance
from sklearn.inspection import permutation_importance

perm_importance = permutation_importance(
    modelo, X_test, y_test,
    n_repeats=10, random_state=42
)

sorted_idx = perm_importance.importances_mean.argsort()[::-1]
plt.boxplot(perm_importance.importances[sorted_idx].T,
            labels=X_test.columns[sorted_idx], vert=False)
\end{verbatim}
\end{frame}

\begin{frame}{Comparación de Métodos}
\small
\begin{table}
\begin{tabular}{llll}
\toprule
\textbf{Método} & \textbf{Scope} & \textbf{Model-agnostic?} & \textbf{Mejor para} \\
\midrule
SHAP & Local/Global & Sí & Explicaciones completas \\
LIME & Local & Sí & Casos individuales \\
PDP & Global & Sí & Efecto de una feature \\
Feature Imp. & Global & Depende & Ranking de features \\
Perm. Imp. & Global & Sí & Importancia real \\
\bottomrule
\end{tabular}
\end{table}

\vspace{0.3cm}

\textbf{Recomendación práctica:}
\begin{enumerate}
\item Feature importance: identificar features relevantes
\item PDP: entender relaciones
\item SHAP: explicaciones completas y robustas
\item LIME: cuando necesites explicar un caso específico a no técnicos
\end{enumerate}
\end{frame}

\section{Resumen}

\begin{frame}{Resumen de la Clase}
\textbf{Ciclo de vida de modelos:}
\begin{itemize}
\item Desarrollo → Deployment → Monitoreo → Reentrenamiento
\item Data drift vs concept drift
\item Versionado y A/B testing
\end{itemize}

\vspace{0.3cm}

\textbf{Serialización:}
\begin{itemize}
\item pickle (estándar), joblib (recomendado para ML), ONNX (interoperable)
\item Guardar pipeline completo + metadata
\end{itemize}

\vspace{0.3cm}

\textbf{Interpretabilidad:}
\begin{itemize}
\item Trade-off con performance
\item SHAP: explicaciones basadas en valores de Shapley
\item LIME: aproximaciones locales
\item PDP: efectos marginales
\item Feature importance: ranking de variables
\end{itemize}
\end{frame}

\begin{frame}{Próxima Clase}
\textbf{Clase 28: Comunicación y Ética}

\vspace{0.3cm}

Veremos:
\begin{itemize}
\item Comunicación efectiva de resultados
\item Reportes reproducibles: Jupyter, R Markdown
\item Visualizaciones para diferentes audiencias
\item Ética en análisis de datos: sesgos, fairness, privacidad
\item Responsabilidad en uso de modelos predictivos
\end{itemize}
\end{frame}

\end{document}
